\subsection{Conservation mechanism ablation}\label{app:conservation}
This experiment compares three ways of obtaining flow conservation on
Sioux Falls under otherwise identical supervision, isolating what the
conservative parameterization of Eq.~\eqref{eq:consparam} contributes
beyond conservation itself:
(i)~\emph{exact} --- the path-share parameterization;
(ii)~\emph{penalty} --- a direct MLP on link flows with a squared
node-balance penalty (weight $1$, normalized by the nodal-throughput
RMS over the training states, floored at the network average);
(iii)~\emph{projection} --- the same direct MLP whose output is
projected each step onto $\{f:Bf=Mq+b\}$ by the differentiable
least-norm map $f\mapsto f-B^{\!\top}(BB^{\!\top})^{\dagger}(Bf-Mq-b)$
(Moore--Penrose inverse; $BB^{\!\top}$ is singular for a connected
network). The direct-MLP width is solved from the parameter-count
equation so that all three models have $714{,}979$ parameters against
the exact model's $715{,}040$. The exact rows reuse the published
five-seed models of Section~\ref{sec:pairresults}; the direct models
are pretrained from scratch on the same data with the same schedule
(Adam $3\times10^{-3}$, batch 256, at most 12{,}000 steps,
best-validation checkpoint), except that their flow loss uses per-link
scaling (necessary for the direct output head to converge), whereas
the published exact models were pretrained on the unscaled loss ---
the one recipe difference in the ablation. Continuation follows the
protocol of Appendix~\ref{app:pairs-v102} unchanged (1{,}200 updates,
shared labels and minibatch sequences, the uniform selection rule
including the pre-trained initial model in the candidate pool), with
flow-only and response supervision for each mechanism. The
exact-mechanism rows reproduce the published medians exactly ($21.18\%$
and $14.19\%$ column error), and the projection output satisfies
$Bf=Mq+b$ to $2\times10^{-11}$.

\begin{table}[pos=htbp]\centering\footnotesize
\caption{Conservation mechanism ablation on the 96 Sioux Falls test
paths (medians of path-level values; active queries and seeds averaged
within path). Residuals are per-query maxima over the path states,
maximized over queries: the node-balance residual
$\max_q\|Bf(q)-Mq-b\|_\infty$ (veh/h); the derivative residual
$\max_{q,j}\|\partial_j[Bf](q)-\mu_j\|_\infty/\|\mu_j\|_2$; the
attribution residual $\max_j\|R_j\|_2/(|\Delta q_j|\|\mu_j\|_2)$
with $R_j$ the integrated flow-attribution residual of
Eq.~\eqref{eq:residualidentity}, computed by trapezoid quadrature of
the learned Jacobian along each test path. Column: attribution column
error; incremental: endpoint flow-change error.}\label{tab:conservation}
\setlength{\tabcolsep}{4pt}
\begin{tabular*}{\tblwidth}{@{\extracolsep{\fill}}llrrrrrr}
\toprule
Mechanism & Method & Flow (\%) & Incremental (\%) & Node resid. & Deriv.\ resid. & Attr.\ resid. & Column (\%)\\
\midrule
Exact & flow & 0.032 & 5.03 & $3\times10^{-11}$ & $9\times10^{-16}$ & $4\times10^{-16}$ & 21.18\\
Exact & response & 0.029 & 3.52 & $3\times10^{-11}$ & $8\times10^{-16}$ & $3\times10^{-16}$ & 14.19\\
\addlinespace
Penalty & flow & 0.040 & 22.76 & 42.1 & 1.16 & 0.66 & 183.7\\
Penalty & response & 0.039 & 22.24 & 42.1 & 1.17 & 0.65 & 177.9\\
\addlinespace
Projection & flow & 0.026 & 14.54 & $2\times10^{-11}$ & $2\times10^{-15}$ & $4\times10^{-15}$ & 104.1\\
Projection & response & 0.026 & 14.42 & $2\times10^{-11}$ & $2\times10^{-15}$ & $4\times10^{-15}$ & 99.4\\
\bottomrule\end{tabular*}\end{table}

Table~\ref{tab:conservation} supports three readings. First, the
penalty model achieves flow accuracy of the same order as the exact
model ($0.040\%$) with a small output residual ($42$ veh/h against
nodal throughputs in the thousands), yet its integrated attribution
residual --- the quantity $R_j$ of Eq.~\eqref{eq:residualidentity}
--- reaches $0.66$ in units of the demand-conservation scale
$|\Delta q_j|\|\mu_j\|_2$, its derivative balance residual is
$1.16$ in the same normalization, and the downstream outcome columns
inherit the damage ($183.7\%$ column error, $323\%$ system
contribution). The chain of Eq.~\eqref{eq:residualidentity} is
exactly what fails: the penalty controls $r(q)$ but not
$\partial_jr$, and the integrated residual $R_j$ --- not the output
residual --- is what the attribution sees. Second, the projection
enforces output conservation exactly and transfers it to derivatives
and to the integrated residual ($4\times10^{-15}$), cutting the
column error to $104.1\%$ --- conservation itself matters --- while
remaining far above the exact parameterization, because projecting an
unconstrained response field onto node balance does not reproduce the
path-share structure that the assignment imposes. Third, response
supervision improves every mechanism (penalty $183.7\to177.9$,
projection $104.1\to99.4$, exact $21.18\to14.19$), with by far the
largest effect inside the exact parameterization: the inductive
benefit of conservation-by-construction and the statistical benefit
of derivative supervision are complementary rather than substitutes.
Because the direct models are pretrained with the per-link scaled
loss while the published exact models use the unscaled loss (the one
recipe difference, stated above), the exact-versus-direct gap should
be read as the joint effect of the parameterization and its training
recipe rather than a pure mechanism isolation; the within-mechanism
comparisons --- the penalty's residual chain above all --- are
unaffected. Under the total-error budget the ordering is unchanged
(normalized $E_{\mathrm{total}}$: $148.8$ and $144.9$ for the penalty,
$46.4$ and $45.6$ for the projection, $11.35$ and $8.68$ for the
exact model).
